\documentclass[12pt, a4paper]{article}
\usepackage[utf8]{inputenc}
\usepackage[ngerman]{babel}
\usepackage[T1]{fontenc}
\usepackage{geometry}
\geometry{a4paper, margin=2.5cm}
\usepackage{hyperref}
\usepackage{graphicx}
\usepackage{listings}
\usepackage{xcolor}
\usepackage[style=alphabetic, backend=biber]{biblatex}
% Code-Listing Style
\lstset{
    backgroundcolor=\color{gray!10},
    basicstyle=\ttfamily\small,
    breaklines=true,
    keywordstyle=\color{blue},
    stringstyle=\color{red},
    commentstyle=\color{green!50!black},
    numbers=left,
    numberstyle=\tiny\color{gray},
    stepnumber=1,
    showstringspaces=false
}

\title{
    \vspace{-2cm}
    \large Hausarbeit im Modul Einführung in Python \\
    \vspace{0.5cm}
    \textbf{\huge Dungeon Autobattler} \\
    \vspace{0.5cm}
    \large Dozent: Dr. Konrad Völkel \\
    Heinrich-Heine-Universität Düsseldorf
}
\author{Felix Weßels \\ Matrikelnummer: 3175987}
\date{\today}

\begin{document}

\maketitle
\thispagestyle{empty}

\newpage
\setcounter{page}{1}

\tableofcontents
\newpage

\section{Einleitung}
\subsection{Motivation}
% Warum dieses Projekt?
\subsection{Zielsetzung}
% Variante der Projektskizze: Was soll die Software am Ende können? (Minimalanforderungen)

\section{Theoretische Grundlagen \& Werkzeuge}
\subsection{Technologien \& Bibliotheken}
% Welche Pakete/Bibliotheken kamen zum Einsatz? (z.B. ruff, mypy, pytest, externe APIs)
\subsection{Mensch vs. KI: Die Rolle von AI Assistance}
% Generelle Vorüberlegungen zum Einsatz von LLMs im Entwicklungsprozess.

\section{Konzept \& Implementierung}
\subsection{Programmstruktur, Kernlogik \& Datenverarbeitung}
% Architektur beschreiben. Vermeiden: Code Zeile für Zeile nacherzählen!
% Stattdessen konzeptionelle Entscheidungen erklären.

\subsection{Warum so?}
% Wahl bestimmter Entwurfsmuster, OOP vs. funktional.
% Welche Alternativen wurden verworfen und warum? (Begründe Entscheidungen!)

\subsection{KI-Konsultation bei Design-Entscheidungen}
% Wurde die KI bei Planung/Architektur gefragt?
% Was wurde übernommen, wo wurde abgewichen (aufgrund menschlicher Logik)?

\subsection{Code-Generierung: Was wurde wie programmiert?}
\textbf{KI-generiert:} % Welche Module/Funktionen/Tests (z.B. Boilerplate, Regex)? Warum war es hier sinnvoll?
\newline
\textbf{Selbst programmiert:} % Kernlogiken, Workarounds, weil KI den Kontext nicht verstand.

\subsection{Prompt-Engineering \& Iteration}
% Wie sah die Zusammenarbeit aus? Wurde Code blind übernommen oder manuell korrigiert (Halluzinationen)?

\section{Ergebnisse}
\subsection{Funktionsnachweis}
% Outputs, Plots oder Screenshots einbinden, um Funktionalität zu belegen.
% Link zum Quelltext: \url{https://github.com/DeinName/DeinRepo} oder Hinweis auf Sciebo-Upload.
\subsection{Beispiel-Durchlauf}
% Zeige, dass auch falsche Nutzereingaben behandelt werden und nicht zum Absturz führen.

\section{Diskussion \& Fazit}
\subsection{Inwiefern wurde das Ziel erreicht?}
\subsection{Herausforderungen \& Abweichung vom Plan}
% Was lief nicht wie geplant? Welche Features aus der Projektskizze wurden weggelassen/geändert?
\subsection{Kritische Würdigung \& Reflexion}
% Was lief gut, wo gab es unerwartete Probleme (Technik oder KI-Zusammenarbeit)? Was würdest du nächstes Mal anders machen?
\subsection{Ausblick}
% Darf subjektiv sein. Welche Features könnten in Zukunft hinzugefügt werden?

\newpage
\printbibliography

\newpage
\appendix
\end{document}